\chapter{Подготовка рабочего места}
\label{app:setup}

Всё, что описано в этом приложении, нужно выполнить один раз перед началом работы
с пособием. Проверки в конце каждого раздела покажут, что установка прошла успешно.

\section{Установка JDK}

Для работы с пособием нужен JDK версии~21 или новее. Обратите внимание: именно
\term{JDK} (\eng{Java Development Kit}), а не JRE~--- в JRE нет компилятора,
и скомпилировать программу не получится.

Наиболее распространённые бесплатные сборки: Eclipse Temurin, Amazon Corretto,
Azul Zulu. Все они собраны из одних и тех же исходников OpenJDK и для задач
пособия равнозначны.

После установки проверьте версию~--- обе команды должны вывести номер~21 или выше:

\begin{lstlisting}[language=bash,caption={Проверка установки JDK},label={lst:app-jdk}]
java -version
javac -version
\end{lstlisting}

Если вторая команда сообщает, что \texttt{javac} не найден, значит установлен JRE
или каталог \texttt{bin} не прописан в переменной \texttt{PATH}.

\section{Переменные среды}

Многие инструменты (Maven, Gradle, среды разработки) ищут Java по переменной
\texttt{JAVA\_HOME}. Она должна указывать на корневой каталог JDK~--- тот, внутри
которого лежит папка \texttt{bin}.

В Windows (PowerShell, от имени администратора~--- для постоянной установки):

\begin{lstlisting}[language=bash,caption={Настройка JAVA\_HOME в Windows},label={lst:app-winenv}]
[Environment]::SetEnvironmentVariable(
    "JAVA_HOME", "C:\Program Files\Java\jdk-21", "Machine")
\end{lstlisting}

В Linux и macOS строки добавляются в \texttt{\textasciitilde/.bashrc} или
\texttt{\textasciitilde/.zshrc}:

\begin{lstlisting}[language=bash,caption={Настройка JAVA\_HOME в Linux и macOS},label={lst:app-nixenv}]
export JAVA_HOME=/usr/lib/jvm/jdk-21
export PATH=$JAVA_HOME/bin:$PATH
\end{lstlisting}

Проверка~--- команда должна напечатать путь к JDK:

\begin{lstlisting}[language=bash,caption={Проверка JAVA\_HOME},label={lst:app-checkenv}]
echo $JAVA_HOME        # Linux, macOS, Git Bash
echo $env:JAVA_HOME    # Windows PowerShell
\end{lstlisting}

\section{Среда разработки}

Подойдёт любая из трёх: IntelliJ~IDEA (редакция Community бесплатна и её
достаточно), Eclipse~IDE for Java Developers или Visual Studio~Code с расширением
Extension Pack for Java. Примеры в пособии не зависят от выбора среды~--- всё
можно собрать и из командной строки.

\section{Проверка: первая программа}

Создайте файл \texttt{Hello.java} с таким содержимым:

\begin{lstlisting}[language=Java,caption={Проверочная программа},label={lst:app-hello}]
public class Hello {
    public static void main(String[] args) {
        System.out.println("Рабочее место готово");
    }
}
\end{lstlisting}

Скомпилируйте и запустите:

\begin{lstlisting}[language=bash,caption={Сборка и запуск},label={lst:app-run}]
javac -encoding UTF-8 Hello.java
java Hello
\end{lstlisting}

Флаг \texttt{-encoding UTF-8} важен: без него компилятор на Windows читает файл
в системной кодировке и русский текст превращается в мусор. Если на экране
появилась строка «Рабочее место готово»~--- всё установлено правильно.

\section{Maven}

Начиная с главы~\ref{ch:06} проекты собираются через Maven. Отдельно его
устанавливать не обязательно: в примерах используется \term{Maven Wrapper}~---
скрипт \texttt{mvnw}, который сам скачивает нужную версию. Но если хочется
поставить Maven глобально, проверка после установки такая:

\begin{lstlisting}[language=bash,caption={Проверка Maven},label={lst:app-mvn}]
mvn -version
\end{lstlisting}

При запуске через wrapper в Windows PowerShell нужен префикс \texttt{.\textbackslash}:
команда \texttt{.\textbackslash mvnw.cmd} вместо \texttt{mvnw.cmd}. Без него
PowerShell не ищет исполняемые файлы в текущем каталоге и сообщает, что команда
не распознана.

\section{Git}

Для главы~\ref{ch:12} понадобится Git. После установки задайте имя и почту~---
они попадают в каждый коммит:

\begin{lstlisting}[language=bash,caption={Первичная настройка Git},label={lst:app-git}]
git --version
git config --global user.name "Иван Иванов"
git config --global user.email "ivanov@example.com"
\end{lstlisting}

Отдельно стоит настроить обработку переводов строк, иначе при работе в смешанной
команде каждый файл будет попадать в историю как «изменённый целиком»:

\begin{lstlisting}[language=bash,caption={Настройка переводов строк},label={lst:app-crlf}]
git config --global core.autocrlf true    # Windows
git config --global core.autocrlf input   # Linux, macOS
\end{lstlisting}

\section{JavaFX}

JavaFX с версии~11 не входит в состав JDK и подключается как обычная зависимость
проекта. Ничего устанавливать заранее не нужно~--- всё необходимое описано
в главе~\ref{ch:09}, где приводится готовый файл \texttt{pom.xml}.
